\title{DeepSeekMath: Pushing the Limits of Mathematical Reasoning in Open Language Models}

\begin{abstract}
Mathematical reasoning remains a formidable task for language models due to its inherent complexity and rigid structure. In this work, we present DeepSeekMath~7B, which extends the pre-training of DeepSeek-Coder-Base-v1.5 7B using 120B tokens containing mathematical content extracted from Common Crawl, combined with text and code samples. DeepSeekMath~7B attains a strong result of 51.7\% on the MATH benchmark at competition-level difficulty, doing so without the need for any external toolkits or ensembling through voting, thereby nearing the benchmark set by Gemini-Ultra and GPT-4. When evaluated with self-consistency across 64 samples, DeepSeekMath~7B reaches a 60.9\% score on MATH. The model's enhanced mathematical reasoning proficiency can be attributed to two principal innovations: firstly, an expertly developed data selection system designed to maximize the utility of freely accessible web datasets; secondly, the introduction of Group Relative Policy Optimization (GRPO), a novel adaptation of Proximal Policy Optimization (PPO), which boosts reasoning performance while simultaneously reducing the memory demands typically associated with PPO.
\end{abstract}

\section{Introduction}

Large language models (LLMs) have fundamentally transformed methods for mathematical reasoning within artificial intelligence, contributing to notable improvements on quantitative reasoning benchmarks \citep{MATH} as well as geometry reasoning assessments \citep{trinh2024solving}. In addition to these developments, LLMs have served as valuable tools for aiding individuals in tackling sophisticated mathematical challenges \citep{tao}. Despite this progress, state-of-the-art systems like GPT-4 \citep{gpt4} and Gemini-Ultra \citep{gemini} remain proprietary, and the open-source alternatives currently available exhibit a marked gap in performance.

This paper presents DeepSeekMath, a specialized language model tailored to mathematics that demonstrates marked improvements over existing open-source models and approaches GPT-4’s performance on established academic tests. Our approach centers around constructing the DeepSeekMath Corpus, a substantial, high-caliber pre-training dataset containing 120B math tokens. The corpus is sourced from the Common Crawl (CC) utilizing a fastText-based classifier \citep{joulin2016fasttext}. For the first round, training relies on examples from OpenWebMath \citep{openwebmath} as the positive set, and incorporates a varied group of other web content as the negative set. After initial training, we use this classifier to uncover further positive instances in the CC, which are then rigorously curated through human labeling. The updated, richer dataset is subsequently used to retrain the classifier, thereby boosting its accuracy. Assessment outcomes demonstrate the corpus’ high quality, evidenced by our baseline model, DeepSeekMath-Base 7B, scoring 64.2\% on GSM8K \citep{gsm8k} and 36.2\% on the advanced MATH dataset \citep{MATH}, surpassing Minerva 540B \citep{minerva}. Furthermore, due to the multilingual nature of the DeepSeekMath Corpus, we observe notable enhancements on Chinese math benchmarks \citep{wei2023cmath,agieval}. We suggest that the methods developed for mathematical data curation presented in this study serve as an initial contribution for further investigations by the research community, with considerable scope for subsequent advancement.

DeepSeekMath-Base is built upon DeepSeek-Coder-Base-v1.5 7B \citep{deepseek-coder}, as our findings suggest that initializing with a code-focused training model yields superior results over leveraging a standard general-purpose LLM. Additionally, math-oriented training not only bolsters mathematical proficiency but also leads to performance gains on MMLU \citep{mmlu} and BBH benchmarks \citep{bbh}, signifying an enhancement in both mathematical and broader reasoning skills.

Following pre-training, we conduct mathematical instruction tuning on DeepSeekMath-Base using datasets for chain-of-thought \citep{cot}, program-of-thought \citep{pot,pal}, and tool-integrated reasoning \citep{tora}. The tuned model, DeepSeekMath-Instruct 7B, outperforms all other 7B models and demonstrates performance on par with instruction-tuned open-source models of 70B scale.

Additionally, we present Group Relative Policy Optimization (GRPO), a reinforcement learning (RL) algorithm variant that builds upon Proximal Policy Optimization (PPO) \citep{schulman2017proximal}. GRPO eliminates the need for a critic model by estimating baselines directly from group scores, which leads to a considerable reduction in the computational requirements during training. Utilizing only a portion of English instruction tuning data, GRPO achieves notable performance gains over the robust DeepSeekMath-Instruct baseline, demonstrating substantial improvements across both in-domain datasets (GSM8K: 82.9\% $\rightarrow$ 88.2\%, MATH: 46.8\% $\rightarrow$ 51.7\%) and out-of-domain tasks (such as CMATH: 84.6\% $\rightarrow$ 88.8\%) in the RL training stage. We further offer a unified framework for interpreting various approaches, such as Rejection Sampling Fine-Tuning (RFT) \citep{yuan2023scaling}, Direct Preference Optimization (DPO) \citep{dpo}, PPO, and GRPO, highlighting that these techniques can be regarded as forms of either direct or streamlined RL. To better understand the critical components of this unified view, we perform comprehensive empirical analyses, including comparisons of online versus offline training, outcome versus process supervision, and single-turn versus iterative RL, among others. Finally, we provide an explanation for the RL-induced performance gains observed in instruction-tuned models and discuss future avenues for developing more effective RL strategies within this consolidated framework.

\subsection{Contributions}
We present scalable pre-training methods for mathematical content, coupled with an in-depth investigation and assessment of reinforcement learning techniques.

\textbf{Math Pre-Training at Scale}

\begin{itemize}[topsep=0pt]
    \item This study demonstrates that the openly available Common Crawl repository includes substantial resources applicable to mathematical tasks. Through a carefully crafted data selection pipeline, we develop the DeepSeekMath Corpus, a large-scale, high-quality dataset comprising 120B tokens drawn from web pages selected for their mathematical relevance. This corpus is nearly 7 times larger than the math web pages utilized by Minerva \citep{minerva} and approximately 9 times the size of the newly introduced OpenWebMath \citep{openwebmath}.

\item DeepSeekMath-Base 7B, our pre-trained base model, performs on par with Minerva 540B \citep{minerva}, suggesting that parameter count alone does not determine mathematical reasoning proficiency. Strong results can also be obtained by training a more compact model on well-curated datasets.

\item We present our observations from conducting math training experiments. Preceding math training with code training enhances models' performance in mathematical problem-solving, whether tools are utilized or not. This provides evidence relevant to the enduring inquiry: \textit{does code training improve reasoning abilities?} Our results suggest that, specifically in the context of mathematical reasoning, it does.

\item While utilizing arXiv papers for training is a widespread approach—particularly among math-focused studies—it yields no significant gains across any of the mathematical benchmarks considered in this study.
\end{itemize}

\textbf{Exploration and Analysis of Reinforcement Learning}

\begin{itemize}[topsep=0pt]
    \item We present Group Relative Policy Optimization (GRPO), a reinforcement learning framework that operates without a critic by estimating the baseline through aggregated group scores. This approach substantially lowers computational demands compared to Proximal Policy Optimization (PPO) while maintaining high effectiveness.
    \item Our experiments show that applying GRPO leads to notable improvements in our instruction-tuned model DeepSeekMath-Instruct when trained solely with instruction-tuning data. Additionally, we observe gains in out-of-domain generalization during the reinforcement learning phase.
    \item We propose a comprehensive framework that integrates diverse methodologies, including RFT, DPO, PPO, and GRPO, under a common perspective. We also perform a thorough investigation featuring a wide range of experimental setups—such as online versus offline training, outcome-based versus process-based supervision, and comparisons between single-turn and iterative reinforcement learning—to closely examine the fundamental factors within this unified paradigm.
    \item Utilizing this unified framework, we analyze the mechanisms underpinning the success of reinforcement learning, and outline a set of promising avenues for further advancing reinforcement learning methodologies for LLMs.
\end{itemize}

\subsection{Summary of Evaluations and Metrics}

\begin{itemize}[topsep=0pt]
    \item \textbf{English and Chinese Mathematical Reasoning}: Our evaluation extensively analyzes model performance across both English and Chinese datasets, incorporating mathematical problems that span from elementary to collegiate levels. For English, the benchmarks utilized are GSM8K \citep{gsm8k}, MATH \citep{MATH}, SAT \citep{llemma}, OCW Courses \citep{minerva}, and MMLU-STEM \citep{mmlu}. For Chinese, we use MGSM-zh \citep{mgsm}, CMATH \citep{wei2023cmath}, Gaokao-MathCloze \citep{agieval}, and Gaokao-MathQA \citep{agieval}. The assessment focuses on two capabilities: producing comprehensive, standalone text-based solutions (without relying on external tools), and solving problems programmatically using Python.

On English evaluation tasks, DeepSeekMath-Base shows performance on par with the proprietary Minerva 540B model \citep{minerva}, and outperforms all existing open-source base models—including Mistral 7B \citep{mistral} and Llemma-34B \citep{llemma}—irrespective of their exposure to mathematical pre-training, frequently by a notable margin. Importantly, DeepSeekMath-Base achieves even better results on Chinese assessment datasets. This advantage likely stems from our divergence from earlier approaches \citep{minerva,llemma}, as we supplement the math pre-training data with high-quality non-English material instead of restricting it to English sources alone. Further enhancements, such as instruction tuning for mathematical reasoning and the incorporation of reinforcement learning, yield DeepSeekMath-Instruct and DeepSeekMath-RL, both of which set a new open-source benchmark by surpassing 50\% accuracy on the MATH dataset—a competition-level challenge.

\item \textbf{Formal Mathematics}: We assess DeepSeekMath-Base on the informal-to-formal theorem proving benchmark introduced by \citep{dsp_proof}, employing the miniF2F dataset \citep{minif2f} and selecting Isabelle \citep{isabelle} as the proof assistant for evaluation. Our results indicate that DeepSeekMath-Base achieves robust few-shot performance in the autoformalization setting.
\item \textbf{Natural Language Understanding, Reasoning, and Code}: To holistically examine the general abilities of DeepSeekMath-Base in comprehension, reasoning, and programming, we conduct experiments on the Massive Multitask Language Understanding (MMLU) suite \citep{mmlu}, which consists of 57 multiple-choice exams spanning a wide range of domains; BIG-Bench Hard (BBH) \citep{bbh}, including 23 tasks focused primarily on complex, multi-step reasoning; as well as HumanEval \citep{codex} and MBPP \citep{mbpp}, two benchmarks commonly used for evaluating coding models. Pre-training on mathematical data proves advantageous for both language understanding and complex reasoning performance.

\section{Math Pre-Training}

\subsection{Data Collection and Decontamination} This section details the methodology used to assemble the DeepSeekMath Corpus utilizing data from Common Crawl. Figure \ref{fig:data_collect} illustrates the iterative workflow developed for efficiently collecting a substantial mathematical dataset from Common Crawl. The process initiates with a seed corpus—comprised of a limited, high-quality math-centric dataset—to anchor subsequent expansion. It is important to emphasize that this method can be similarly extended to other fields, including programming.

We begin by selecting OpenWebMath \citep{openwebmath}, an established repository of high-quality mathematical web material, as our foundational seed dataset. Leveraging this dataset, we construct and train a fastText model \citep{joulin2016fasttext} aimed at retrieving additional web pages with characteristics similar to those in OpenWebMath. For training purposes, we randomly designate 500,000 entries from the seed dataset as positive samples and an equal number of Common Crawl web pages as negative samples. Training is carried out using an open-source implementation\footnote{\url{https://fasttext.cc}}, with the following hyperparameters: vector size of 256, a learning rate set to 0.1, a maximum word n-gram length of 3, a minimum word frequency threshold of 3, and 3 epochs. To downsize the Common Crawl dataset, both URL-based and near-duplication removal are applied, producing a refined corpus of 40B HTML documents. The trained fastText model is subsequently employed to retrieve mathematical pages from this deduplicated set. In order to exclude mathematical pages of substandard quality, we rank the retrieved documents by their model-assigned scores, retaining only those with the highest ranks. The efficacy of this filtering and ranking procedure is further examined by conducting pre-training experiments using token subsets of varying sizes—specifically, 40B, 80B, 120B, and 160B tokens from the highest-scoring pages. For the first iteration, the top 40B tokens are selected for preservation.

Following the initial round of data collection, a significant portion of mathematical web pages remain uncaptured. This is primarily attributed to the limited diversity present in the set of positive examples used to train the fastText model. To address this, we expand our pool of seed data by identifying and incorporating additional mathematical web domains. Our method involves segmenting the complete Common Crawl dataset according to unique domains—each defined by web pages with a shared base URL. We then compute, for every domain, the proportion of web pages acquired during the first collection phase. Domains with more than 10\% of their web pages collected are tagged as being math-oriented (for instance, \url{mathoverflow.net}).

We proceed by manually annotating URLs within these math-centric domains that are indicative of mathematical content (e.g., \url{mathoverflow.net/questions}). Any web pages connected to these URLs but not previously gathered are subsequently incorporated into the seed corpus. By augmenting the seed corpus in this manner, we can assemble a broader set of positive training examples, thus enabling the fastText model to identify a greater quantity of mathematical data in future rounds. This iterative process continues, and after four cycles, we assemble a dataset consisting of 35.5M mathematical web pages, cumulatively amounting to 120B tokens. Analysis after the fourth iteration shows that approximately 98\% of mathematical web pages were already acquired by the third iteration, prompting us to halt further data collection.

To mitigate the risk of benchmark data leakage, we adopt the filtering strategy described in \cite{deepseek-coder}, excluding web content that includes questions or answers sourced from English math benchmarks like GSM8K \citep{gsm8k} and MATH \citep{MATH} as well as Chinese benchmarks such as CMATH \citep{wei2023cmath} and AGIEval \citep{agieval}. Specifically, our approach removes any segment of text from the math training set if it contains a contiguous 10-gram phrase that perfectly matches any sub-sequence from the benchmark datasets. For benchmark content that falls below the 10-gram threshold but consists of at least 3 grams, we rely on an exact matching process to filter out potentially contaminated web material.

\subsection{Validating the Quality of the DeepSeekMath~Corpus}

We run pre-training experiments to investigate how the DeepSeekMath~Corpus is compared with the recently released math-training corpora:

\begin{itemize}[topsep=0pt]
    \item \textbf{MathPile} \citep{mathpile}: This dataset comprises 8.9B tokens drawn from a variety of sources, including textbooks, Wikipedia, ProofWiki, CommonCrawl, StackExchange, and arXiv, with the latter providing more than 85\% of the content.
    \item \textbf{OpenWebMath} \citep{openwebmath}: Extracted from CommonCrawl by isolating mathematical material, this corpus amounts to 13.6B tokens.
    \item \textbf{Proof-Pile-2} \citep{llemma}: This resource merges several mathematical datasets—OpenWebMath, AlgebraicStack (offering 10.3B tokens of math-related code), and arXiv scholarly articles (amounting to 28.0B tokens). For experiments involving Proof-Pile-2, we adopt the 2:4:1 arXiv:Web:Code ratio specified by \cite{llemma}.
\end{itemize}

\subsubsection{Training Setting}
\label{sec:quality-policy}
Mathematical fine-tuning is performed on a general-purpose language model containing 1.3B parameters, architecturally consistent with DeepSeek LLMs \citep{deepseek-llm}, and referred to as DeepSeek-LLM 1.3B. Separate models are fine-tuned on each respective mathematical dataset for a total of 150B tokens. The HAI-LLM framework \citep{haillm}, known for its efficiency and minimal overhead, is employed for all training experiments. Consistent with DeepSeek LLMs' methodologies, optimization is carried out using AdamW \citep{adamW}, setting $\beta_1=0.9$, $\beta_2=0.95$, and $\mathrm{weight\_decay}=0.1$. The learning rate schedule follows a staged approach, increasing linearly to its peak after 2,000 warmup steps, reducing to 31.6\% at the 80\% mark of training, and further tapering to 10.0\% of the maximum after 90\% of training is completed. The maximum learning rate is fixed at 5.3e-4. Each batch contains 4 million tokens, and model context length is set to 4,000 tokens.

\subsubsection{Evaluation Results}

\textbf{The DeepSeekMath~Corpus is of high quality, covers multilingual mathematical content, and is the largest in size.}

\begin{itemize}[topsep=0pt]
    \item \textbf{High-quality}:
    We assess downstream capabilities across 8 different mathematical benchmarks utilizing few-shot chain-of-thought prompting \cite{cot}. As summarized in Table \ref{tab:corpora_comparison}, the model trained with the DeepSeekMath~Corpus achieves a distinct performance advantage. Furthermore, Figure \ref{fig:corpora_comparisons} illustrates that, at the 50B token mark (corresponding to one epoch on Proof-Pile-2), the DeepSeekMath Corpus-trained model consistently surpasses Proof-Pile-2, suggesting a superior average quality within the DeepSeekMath Corpus. 
    \item \textbf{Multilingual}:
    The DeepSeekMath~Corpus includes mathematical texts in various languages, with English and Chinese constituting the largest segments. According to Table \ref{tab:corpora_comparison}, employing the DeepSeekMath~Corpus during training produces marked gains in mathematical reasoning tasks for both English and Chinese content. Conversely, other available mathematical datasets, which focus predominantly on English, deliver minimal or even negative transfer to Chinese mathematical reasoning performance.
    \item \textbf{Large-scale}:
    In terms of scale, the DeepSeekMath~Corpus substantially exceeds other mathematical datasets. Figure \ref{fig:corpora_comparisons} demonstrates that DeepSeek-LLM 1.3B, when trained on the DeepSeekMath~Corpus, not only displays a more accelerated learning trajectory but also sustains improvements over time. By comparison, other corpora have much smaller sizes and are cycled through several times during training, which leads to models reaching peak performance much earlier.
\end{itemize}

\subsection{Training and Evaluating DeepSeekMath-Base 7B}

This section presents DeepSeekMath-Base 7B, a foundational model that exhibits notable proficiency in mathematical reasoning. The model builds upon DeepSeek-Coder-Base-v1.5 7B \citep{deepseek-coder} as its starting point and undergoes training across 500B tokens. The data utilized during training comprises 56\% from the DeepSeekMath Corpus, 4\% from AlgebraicStack, 10\% originating from arXiv, 20\% from Github code, with the final 10\% consisting of natural language content drawn from Common Crawl in both English and Chinese. Our training procedure largely aligns with that described in Section \ref{sec:quality-policy}; however, we adjust the peak learning rate to 4.2e-4 and employ a batch size of 10M tokens.

In this work, we thoroughly evaluate DeepSeekMath-Base 7B's mathematical proficiency by examining its capacity to independently generate complete mathematical solutions, utilize external tools for problem-solving, and engage in formal theorem proving. Additionally, our analysis extends beyond mathematical tasks to encompass a broader overview of the base model's general abilities, assessing its competence in natural language understanding, reasoning, and programming.

\paragraph{Mathematical Problem Solving with Step-by-Step Reasoning}

We assess DeepSeekMath-Base’s capabilities in mathematical problem-solving through few-shot chain-of-thought prompting \citep{cot}, employing eight evaluation benchmarks presented in both English and Chinese. These benchmarks span tasks involving quantitative reasoning—such as GSM8K \citep{gsm8k}, MATH \citep{MATH}, and CMATH \citep{wei2023cmath}—as well as multiple-choice formats, including MMLU-STEM \citep{mmlu} and Gaokao-MathQA \citep{agieval}. Collectively, they reflect a wide spectrum of mathematical subjects, ranging from elementary up to college-level challenges.

As indicated in Table \ref{tab:base_math_cot}, DeepSeekMath-Base 7B achieves the highest scores across all eight benchmarks when compared to other open-source base models, such as the widely adopted Mistral 7B \citep{mistral} and the more recent Llemma 34B \citep{llemma}, the latter being fine-tuned on Proof-Pile-2 \citep{llemma}. Importantly, DeepSeekMath-Base surpasses all open-source base model contenders on the MATH dataset—targeting competition-level tasks—by more than 10\% in absolute terms. Additionally, it outperforms the proprietary Minerva 540B \citep{minerva}, a model 77 times larger which is based on PaLM \citep{palm} and further specialized with mathematical content.

\paragraph{Mathematical Problem Solving with Tool Use} We assess mathematical reasoning supported by program execution on the GSM8K and MATH datasets, employing few-shot program-of-thought prompting as described by \citet{pot,pal}. In this setup, models approach each task by generating a Python script that may use packages like \textit{math} and \textit{sympy} for performing complex calculations. The output derived from running the code is considered the solution to the problem. As detailed in Table \ref{tab:base_math_pot_proof}, DeepSeekMath-Base 7B demonstrates superior performance compared to the earlier leading model Llemma 34B.

\paragraph{Formal Mathematics}

Automating formal proofs plays a critical role in promoting the correctness and dependability of mathematical reasoning, as well as improving workflow efficiency—a topic that has garnered growing interest recently. In this work, we assess DeepSeekMath-Base 7B using the informal-to-formal proof task outlined in \citep{dsp_proof}, which involves constructing a formal proof given an informal statement, its formal representation, and an informal derivation. Our evaluation utilizes the miniF2F dataset \citep{minif2f}, recognized as a benchmark for Olympiad-level formal mathematics, where we produce Isabelle-based formal proofs for each problem through few-shot prompting. Consistent with the methodology from \cite{dsp_proof}, the model is first tasked with generating a proof outline, and then the Sledgehammer automated prover \citep{sledgehammer} is used to automatically supply any missing inferential steps. As summarized in Table \ref{tab:base_math_pot_proof}, DeepSeekMath-Base 7B achieves notable results in the domain of proof autoformalization.

\paragraph{Natural Language Understanding, Reasoning, and Code}

We assess the capabilities of models in natural language understanding using MMLU \citep{mmlu}, in reasoning with BBH \citep{bbh}, and in coding tasks via HumanEval \citep{codex} and MBPP \citep{mbpp}. According to the results presented in Table \ref{tab:understanding_reasoning_code}, DeepSeekMath-Base 7B demonstrates substantial improvement on MMLU and BBH in comparison to its predecessor, DeepSeek-Coder-Base-v1.5 \citep{deepseek-coder}, highlighting the beneficial effects of mathematical training on both language comprehension and reasoning abilities. Furthermore, through the inclusion of code tokens in continued training, DeepSeekMath-Base 7B preserves the strong coding benchmark performance characteristic of DeepSeek-Coder-Base-v1.5. Taken together, DeepSeekMath-Base 7B achieves markedly higher scores than the general-purpose Mistral 7B model \citep{mistral} across all three reasoning and coding evaluations.

\section{Supervised Fine-Tuning}

\subsection{SFT Data Curation}

We develop a dataset for mathematical instruction tuning that encompasses both English and Chinese questions drawn from a variety of mathematical disciplines and complexity tiers. Each question is matched with solutions using chain-of-thought (CoT) \citep{cot}, program-of-thought (PoT) \citep{pot,pal}, and tool-integrated reasoning methodologies \citep{tora}. The resulting corpus consists of 776K training samples.

\begin{itemize}[topsep=0pt]
    \item \textbf{English mathematical datasets}: We provide annotated tool-integrated solutions for GSM8K and MATH datasets, and incorporate a selected portion of MathInstruct \citep{MathInstruct}, together with the training partition of Lila-OOD \citep{lila}, in which the problems are solved using either CoT or PoT reasoning. The English dataset comprises a broad spectrum of mathematical disciplines, including but not limited to algebra, probability, number theory, calculus, and geometry.
    \item \textbf{Chinese mathematical datasets}: We aggregate K-12 mathematics problems in Chinese, covering 76 different sub-topics such as linear equations, with each problem annotated using both CoT-based and tool-integrated solution methodologies.
\end{itemize}

\subsection{Training and Evaluating DeepSeekMath-Instruct 7B}

This section presents DeepSeekMath-Instruct 7B, a model further refined through mathematical instruction tuning built upon DeepSeekMath-Base. For training, examples are combined in random order until the combined token sequence approaches a maximum context length of 4K tokens. The model undergoes 500 training steps using a batch size of 256 and a fixed learning rate set at 5e-5.

We evaluate models' mathematical performance both without and with tool use, on 4 quantitative reasoning benchmarks in English and Chinese. We benchmark our model against the leading models of the time:

\begin{itemize}[topsep=0pt]
    \item \textbf{Closed-source models} encompass the following: (1) the GPT series, with GPT-4 \citep{gpt4} and GPT-4 Code Interpreter~\footnote{\url{https://openai.com/blog/chatgpt-plugins\#\#code-interpreter}} representing the most advanced variants, (2) Gemini Ultra and Pro \citep{gemini}, (3) Inflection-2 \citep{inflection-2}, (4) Grok-1~\footnote{\url{https://x.ai/model-card}}, in addition to models recently introduced by Chinese enterprises, such as (5) Baichuan-3~\footnote{\url{https://www.baichuan-ai.com}}, (6) and the latest iteration, GLM-4~\footnote{\url{https://open.bigmodel.cn/dev/api\#glm-4}}
\end{itemize}

    from the GLM family \citep{glm}.

These models are designed for broad applicability and have typically been subjected to multiple alignment stages. 
    \item \textbf{Open-source models} encompass: general-purpose systems such as (1) DeepSeek-LLM-Chat 67B \citep{deepseek-llm}, (2) Qwen 72B \citep{qwen}, (3) SeaLLM-v2 7B \citep{seallm}, and (4) ChatGLM3 6B \citep{chatglm3}. The category also features models with particular strengths in mathematics, including: (5) InternLM2-Math 20B~\footnote{\url{https://github.com/InternLM/InternLM-Math}}, an extension of InternLM2 that has undergone specialized mathematical training as well as subsequent instruction tuning; (6) Math-Shepherd-Mistral 7B, which utilizes PPO training \citep{schulman2017proximal} on the Mistral 7B \citep{mistral} platform with a reward model based on process supervision; (7) the WizardMath series \citep{wizardmath}, which enhances mathematical reasoning abilities in Mistral 7B and Llama-2 70B \citep{llama2} by leveraging evolve-instruct (an instruction tuning variant employing AI-generated instructions) and PPO, drawing most of their training data from GSM8K and MATH datasets; (8) MetaMath 70B \citep{metamath}, which represents Llama-2 70B that has been fine-tuned on enriched GSM8K and MATH resources; (9) ToRA 34B \cite{tora}, a CodeLlama 34B variant trained specifically for tool-based mathematical reasoning; and (10) MAmmoTH 70B \citep{MathInstruct}, which is based on Llama-2 70B with instruction tuning carried out on the MathInstruct dataset.
\end{itemize}

Table \ref{tab:sft_rl_math} illustrates that, in the setting where external tool use is prohibited, DeepSeekMath-Instruct 7B excels in step-by-step mathematical reasoning. Remarkably, on the competition-tier MATH dataset, our model achieves at least a 9\% absolute improvement over all open-source alternatives and most closed-source systems (including Inflection-2 and Gemini Pro), even outperforming models with substantially larger parameter counts (e.g., Qwen 72B) or those specifically tuned with math-oriented reinforcement learning such as WizardMath-v1.1 7B. Although DeepSeekMath-Instruct achieves comparable results to prominent Chinese proprietary systems GLM-4 and Baichuan-3 on this benchmark, its performance still lags behind GPT-4 and Gemini Ultra.

In an evaluation framework permitting models to combine natural language reasoning with tool-assisted programmatic approaches for solving problems, DeepSeekMath-Instruct 7B nearly attains a 60\% accuracy rate on MATH, outperforming all previously released open-source models. Across other benchmarks, the model maintains performance on par with DeepSeek-LLM-Chat 67B—the former state-of-the-art, despite its tenfold larger parameter count.  
\section{Reinforcement Learning}

\subsection{Group Relative Policy Optimization} Following Supervised Fine-Tuning (SFT), reinforcement learning (RL) has demonstrated its value in enhancing LLMs' capacity for mathematical reasoning \citep{wang2023math,wizardmath}. Here, we present our novel RL method, Group Relative Policy Optimization (GRPO), which offers both efficiency and efficacy.

\subsubsection{From PPO to GRPO} Proximal Policy Optimization (PPO) \citep{schulman2017proximal} is an actor-critic RL algorithm that is widely used in the RL fine-tuning stage of LLMs \citep{ouyang2022training}. In particular, it optimizes LLMs by maximizing the following surrogate objective:

\begin{equation}
\footnotesize
    \mathcal{J}_{PPO}(\theta) = \mathbb{E}{[q \sim P(Q),\ o \sim \pi_{\theta_{old}}(O|q)]} \frac{1}{|o|} \sum_{t=1}^{|o|} \min \left[ \frac{\pi_\theta(o_{t} | q, o_{<t})}{\pi_{\theta_{old}}(o_{t} | q, o_{<t})} A_{t},\ \text{clip}\left(\frac{\pi_\theta(o_{t} | q, o_{<t})}{\pi_{\theta_{old}}(o_{t} | q, o_{<t})},\ 1 - \varepsilon,\ 1 + \varepsilon \right) A_{t} \right],
\end{equation}

The objective function $\mathcal{J}_{PPO}(\theta)$ leverages expectations over $q$ sampled from $P(Q)$ and $o$ from the prior policy $\pi_{\theta_{old}}(O|q)$. For each token in the output sequence, the loss computes the average across sequence length $|o|$. At each step $t$, it takes the minimum between the ratio of likelihoods under current and old policies scaled by $A_t$, and a clipped variant of this ratio (within $1 \pm \varepsilon$) also multiplied by $A_t$, thus preventing large policy updates and ensuring training stability.

In this context, $\pi_{\theta}$ denotes the current policy model, while $\pi_{\theta_{old}}$ refers to the previous version. The variables $q$ and $o$ represent questions and their corresponding outputs, sampled respectively from the question dataset and generated by the earlier policy $\pi_{\theta_{old}}$. The hyper-parameter $\varepsilon$ is used for clipping within the PPO framework, serving to promote more stable training dynamics. The advantage term $A_t$ is estimated via Generalized Advantage Estimation (GAE) \citep{gae}, utilizing the reward sequence $\{r_{\ge t}\}$ and a parameterized value function $V_{\psi}$. Accordingly, in PPO, the optimization involves learning a value function in tandem with the policy. To address the potential over-optimization of the reward model, it is standard practice to introduce a KL penalty term computed per token, measuring divergence from a reference model and incorporated into the reward signal for each token \citep{ouyang2022training}; that is,

\begin{equation}
    r_{t} = r_\varphi(q, o_{\le t}) - \beta \log\frac{\pi_{\theta}(o_{t}|q, o_{<t})}{\pi_{ref}(o_{t}|q, o_{<t})},
\label{eq:PPO-reward}
\end{equation}

Here, $r_\varphi$ represents the reward model, $\pi_{ref}$ denotes the reference model—commonly set as the initial SFT model—and $\beta$ indicates the weight assigned to the KL penalty term.

As the value function employed in PPO is typically another model of comparable size as the policy model, it brings a substantial memory and computational burden. Additionally, during RL training, the value function is treated as a baseline in the calculation of the advantage for variance reduction. While in the LLM context, usually only the last token is assigned a reward score by the reward model, which may complicate the training of a value function that is accurate at each token. To address this, as shown in Figure \ref{fig:grpo}, we propose Group Relative Policy Optimization (GRPO), which obviates the need for additional value function approximation as in PPO, and instead uses the average reward of multiple sampled outputs, produced in response to the same question, as the baseline. More specifically, for each question $q$, GRPO samples a group of outputs $\{o_1, o_2, \cdots, o_G\}$  from the old policy  $\pi_{\theta_{old}}$  and then optimizes the policy model by maximizing the following objective:

\begin{equation}
\footnotesize
\begin{split}
    \mathcal{J}_{GRPO}(\theta) &= \mathbb{E}{[q \sim P(Q), \{o_i\}_{i=1}^G \sim \pi_{\theta_{old}}(O|q)]}  \\
    & \frac{1}{G}\sum_{i=1}^G\frac{1}{|o_i|} \sum_{t=1}^{|o_i|} \Bigg\{ \min \Bigg[ \frac{\pi_\theta(o_{i,t} | q, o_{i,<t})}{\pi_{\theta_{old}}(o_{i,t} | q, o_{i,<t})} \hat{A}_{i,t},\ \text{clip} \Bigg( \frac{\pi_\theta(o_{i,t} | q, o_{i,<t})}{\pi_{\theta_{old}}(o_{i,t} | q, o_{i,<t})},\ 1 - \varepsilon,\ 1 + \varepsilon \Bigg) \hat{A}_{i,t} \Bigg] - \beta \mathbb{D}_{KL}\Big[\pi_{\theta}\ ||\ \pi_{ref}\Big]\Bigg\} ,
\end{split}
\label{eq:GRPO-obj}
\end{equation}

where $\varepsilon$ and $\beta$ are hyper-parameters, and $\hat{A}_{i,t}$  is the advantage calculated based on relative rewards of the outputs inside each group only, which will be detailed in the following subsections. The group relative way that GRPO leverages to calculate the advantages, aligns well with the comparative nature of rewards models,  as reward models are typically trained on datasets of comparisons between outputs on the same question. Also note that, instead of adding KL penalty in the reward, GRPO regularizes by directly adding the KL divergence between the trained policy and the reference policy to the loss, avoiding complicating the calculation of  $\hat{A}_{i,t}$. And different from the KL penalty term used in (\ref{eq:PPO-reward}), we estimate the KL divergence with the following unbiased estimator \citep{kl_approx}:

\begin{equation}
\small
    \mathbb{D}_{KL}\left[\pi_{\theta} || \pi_{ref}\right] = \frac{\pi_{ref}(o_{i,t}|q,o_{i,<t})}{\pi_{\theta}(o_{i,t}|q,o_{i,<t})}- \log\frac{\pi_{ref}(o_{i,t}|q,o_{i,<t})}{\pi_{\theta}(o_{i,t}|q,o_{i,<t})} - 1,
\end{equation}

which is guaranteed to be positive.

\begin{algorithm}[t]
  \small
  \caption{Iterative Group Relative Policy Optimization}
  \textbf{Input} initial policy model $\pi_{\theta_{\text{init}}}$; reward models $r_\varphi$; task prompts $\mathcal{D}$; 
  hyperparameters $\varepsilon$, $\beta$, $\mu$
  \begin{algorithmic}[1]
    \State Set policy model $\pi_\theta \leftarrow \pi_{\theta_{\text{init}}}$
    \For{iteration = 1, \dots, I}
       \State Assign current policy as reference model: $\pi_{ref} \leftarrow \pi_{\theta}$
      \For{step = 1, \dots, M}
      \State Select a batch $\mathcal{D}_b$ sampled from $\mathcal{D}$
      \State Save present policy as $\pi_{\theta_{old}} \leftarrow \pi_{\theta}$ 
      \State For each prompt $q \in \mathcal{D}_b$, generate $G$ outputs $\{o_i\}_{i=1}^G \sim \pi_{\theta_{old}} (\cdot \mid q)$ 
      \State For every generated output $o_i$, obtain rewards $\{r_i\}_{i=1}^{G}$ via $r_\varphi$
      \State Derive the estimated group-relative advantage $\hat{A}_{i,t}$ for each token $t$ in $o_i$
      \For{GRPO iteration = 1, \dots, $\mu$}
        \State Optimize the policy model $\pi_{\theta}$ by maximizing the group-relative objective from Equation \ref{eq:GC-GRPO}
      \EndFor
    \EndFor \State Continue updating $r_\varphi$ utilizing an ongoing replay-based training strategy.
    \EndFor 
  \end{algorithmic}
  \textbf{Output} $\pi_\theta$
  \label{alg:iter-grpo}
\end{algorithm}

\subsubsection{Outcome Supervision RL with GRPO} Specifically, given a question $q$, a collection of outputs $\{o_1, o_2, \cdots, o_G\}$ is generated by sampling from the prior policy $\pi_{\theta_{old}}$. Each output is assigned a score by a reward model, resulting in a set of $G$ rewards $\mathbf{r}=\{r_1, r_2, \cdots, r_G\}$. These reward values are standardized by removing the group mean and dividing by the group standard deviation. Outcome supervision applies this normalized score as the reward at the conclusion of each generated output $o_i$, assigning the normalized reward as the advantage $\hat{A}_{i, t}$ for every token within the sequence, so that $\hat{A}_{i, t} = \widetilde{r}_i = \frac{r_i- {\rm mean}(\mathbf{r})}{{\rm std}(\mathbf{r})}$. Finally, the policy parameters are updated by maximizing the objective presented in equation (\ref{eq:GRPO-obj}).

\subsubsection{Process Supervision RL with GRPO}

Relying solely on outcome supervision—which assigns a reward only after generating the entire output—often fails to provide adequate and effective feedback for guiding the policy through intricate mathematical tasks. In line with the approach taken by \cite{wang2023math}, we further investigate process supervision, which allocates rewards at the conclusion of each intermediate reasoning step. More precisely, considering a question $q$ and a set of $G$ sampled outputs $\{o_1, o_2, \cdots, o_G\}$, we utilize a process reward model that assesses each step, resulting in a collection of rewards: $\mathbf{R} = \{ \{r_1^{index(1)},\cdots,r_1^{index(K_1)}\}, \cdots,  \{r_G^{index(1)},\cdots,r_G^{index(K_G)}\} \}$. Here, $index(j)$ specifies the token index corresponding to the end of the $j$-th step, and $K_i$ denotes how many steps are present in the $i$-th output. To ensure uniformity and comparability, these reward values undergo normalization using both the mean and standard deviation: $\widetilde{r}_i^{index(j)} = \frac{r_i^{index(j)} - {\rm mean(\mathbf{R})}}{{\rm std(\mathbf{R})}}$.

Following normalization, the process supervision strategy computes the advantage for each token as the accumulated sum of normalized rewards starting from the current token onward: $\hat{A}_{i, t} = \sum_{index(j) \ge t} \widetilde{r}_i^{index(j)}$. The policy is then updated by maximizing the objective function outlined in equation (\ref{eq:GRPO-obj}).

\subsubsection{Iterative RL with GRPO} During reinforcement learning, the previous reward model may become inadequate for guiding the updated policy model. To address this, we investigate an iterative RL strategy using GRPO. As illustrated in Algorithm \ref{alg:iter-grpo}, this iterative GRPO framework involves generating fresh reward model training data from samples produced by the current policy model. The reward model is then updated in an ongoing manner, utilizing a replay buffer that preserves 10\% of past samples. Subsequently, the updated policy model is adopted as the new reference model, and training of the policy model proceeds under supervision from the newly updated reward model.

\subsection{Training and Evaluating DeepSeekMath-RL}

We perform reinforcement learning utilizing the DeepSeekMath-Instruct 7B model. The RL phase is trained on chain-of-thought-format questions pertaining specifically to GSM8K and MATH, extracted from the SFT dataset, totaling approximately 144,000 items. We deliberately omit other SFT-derived questions to analyze how RL influences benchmark performance when those benchmarks are absent from the RL training data. Our reward model's training data is assembled in accordance with the methodology outlined by \citep{wang2023math}. The initial reward model is trained starting from DeepSeekMath-Base 7B, employing a learning rate of $2 \times 10^{-5}$. When training using GRPO, the policy model utilizes a learning rate of $1 \times 10^{-6}$, with the KL penalty coefficient fixed at 0.04. For each individual question, we generate 64 sampled responses. The maximum sequence length is capped at 1024 tokens, and each training batch also contains 1024 instances. After each exploration round, the policy model undergoes a single parameter update. The assessment of DeepSeekMath-RL 7B across benchmarks mirrors the procedures applied to DeepSeekMath-Instruct 7B. Specifically for DeepSeekMath-RL 7B, benchmarks GSM8K and MATH (when incorporating chain-of-thought reasoning) are treated as in-domain, while all other benchmarks fall into the out-of-domain category.

Table \ref{tab:sft_rl_math} presents results comparing open- and closed-source models that employ both chain-of-thought and tool-augmented reasoning approaches on benchmarks in English and Chinese. The analysis reveals the following: 1) DeepSeekMath-RL 7B achieves accuracy rates of 88.2\% on GSM8K and 51.7\% on MATH using chain-of-thought reasoning, outperforming every open-source competitor in the 7B to 70B parameter range, as well as most closed-source alternatives. 2) Notably, DeepSeekMath-RL 7B’s training involved solely chain-of-thought-style instruction tuning data from GSM8K and MATH datasets, initiated from DeepSeekMath-Instruct 7B. Even with this limited training dataset, DeepSeekMath-RL 7B consistently exceeds the performance of DeepSeekMath-Instruct 7B across all measured criteria, highlighting the benefits conferred by reinforcement learning.

\section{Discussion}
Here, we present an overview of our results from both the pre-training phase and the reinforcement learning (RL) experiments.

\subsection{Lessons Learnt in Pre-Training}

To begin, we discuss our pre-training experience. Except where explicitly indicated, we follow the training configurations detailed in Section \ref{sec:quality-policy}. Notably, in this section, any mention of the DeepSeekMath Corpus pertains to an 89B-token dataset obtained during the second round of data collection.

\subsubsection{Code Training Benefits Mathematical Reasoning}

An often-cited but as yet unsubstantiated assumption is that training on code enhances reasoning skills. In this work, we seek to partially address this claim, focusing on mathematical contexts: code training enhances the mathematical reasoning capabilities of models, irrespective of whether tool use is involved.

To study how code training affects mathematical reasoning, we experimented with the following two-stage training and one-stage training settings:

\textbf{Two-Stage Training}

\begin{itemize}[topsep=0pt]
    \item \textbf{Code Training for 400B Tokens $\rightarrow$ Math Training for 150B Tokens}: The DeepSeek-LLM 1.3B model is initially pretrained on 400B tokens sourced from code, and subsequently fine-tuned on 150B math-specific tokens;
    \item \textbf{General Training for 400B Tokens $\rightarrow$ Math Training for 150B Tokens}: For comparison, we perform an alternative experiment in which the first stage uses 400B tokens drawn from a general-domain corpus assembled by DeepSeek-AI, rather than code tokens, to evaluate whether code-based pretraining brings distinct improvements to mathematical reasoning tasks over general-domain pretraining.
\end{itemize}

\textbf{One-Stage Training}

\begin{itemize}[topsep=0pt]
    \item \textbf{Math Training with 150B Tokens}: DeepSeek-LLM 1.3B is trained using a dataset composed of 150B math-specific tokens;
    \item \textbf{Joint Training with 400B Code Tokens and 150B Math Tokens}: Sequential math training after code training leads to diminished coding ability. Thus, we examine if combining math and code tokens in a unified training phase can foster mathematical reasoning while simultaneously addressing issues like catastrophic forgetting.
\end{itemize}

\paragraph{Results}
The downstream results corresponding to various training configurations are presented in Table \ref{tab:code-math} and Table \ref{tab:code-math-general-eval}.

Training on code brings measurable gains to program-aided mathematical reasoning across both the two-stage and one-stage training frameworks. As illustrated in Table \ref{tab:code-math}, within the two-stage approach, focusing solely on code training already markedly boosts performance on GSM8K and MATH tasks when utilizing Python. Adding mathematical reasoning tasks during the second stage leads to additional performance gains. Notably, in the one-stage regime, interleaving code tokens and math tokens during training proves effective for circumventing the catastrophic forgetting observed in sequential two-stage setups, while concurrently enhancing both general coding ability (Table \ref{tab:code-math-general-eval}) and program-assisted mathematical problem solving (Table \ref{tab:code-math}).

Training with code additionally enhances mathematical reasoning skills even in the absence of external tools. In the two-stage training framework, the preliminary phase that focuses on code already produces noticeable gains. This approach further accelerates and enhances the effectiveness of the following mathematical training phase, culminating in superior performance. In contrast, merging code tokens and math tokens in a single-stage training regimen tends to undermine mathematical reasoning when tools are not available. A possible explanation is that DeepSeek-LLM 1.3B, being relatively small in scale, may not possess sufficient capacity to comprehensively integrate both code and mathematical information at the same time.

\subsubsection{ArXiv Papers Seem Ineffective in Improving Mathematical Reasoning}

ArXiv papers are commonly included as a component of math pre-training data \citep{minerva,gpt-f,llemma,mathpile}. However, detailed analysis regarding their impact on mathematical reasoning has not been extensively conducted. Perhaps counter-intuitively, according to our experiments, arXiv papers seem ineffective in improving mathematical reasoning. We experiment with models of different sizes, including DeepSeek-LLM 1.3B and DeepSeek-Coder-Base-v1.5 7B \citep{deepseek-coder}, using arXiv corpora that underwent varied processing pipelines:

\begin{itemize}[topsep=0pt]
    \item \textbf{MathPile} \citep{mathpile}: a dataset containing 8.9 billion tokens curated through specific cleaning and filtering heuristics, with scientific papers from arXiv comprising more than 85\% of its content;
    \item \textbf{ArXiv-RedPajama} \citep{redpajama}: a 28.0 billion-token compilation of all arXiv LaTeX sources after eliminating preambles, comments, macros, and bibliography sections.
\end{itemize}

For our experimental setup, we individually trained DeepSeek-LLM 1.3B using 150B tokens and DeepSeek-Coder-Base-v1.5 7B with 40B tokens, each on a dedicated arXiv corpus. The results indicate that training on arXiv papers does not meaningfully enhance mathematical reasoning. Models trained exclusively with arXiv data exhibit little to no improvement—and, in some cases, a decline—when evaluated on a range of mathematical benchmarks of varying difficulty in this work. These assessments encompass quantitative reasoning datasets such as GSM8K and MATH (see Table \ref{tab:arxiv-cot}), multiple-choice evaluations like MMLU-STEM (Table \ref{tab:arxiv-cot}), as well as formal mathematics benchmarks including miniF2F (Table \ref{tab:arxiv_atp}).

However, this conclusion has its limitations and should be taken with a grain of salt. We have not yet studied:

\begin{itemize}[topsep=0pt]
    \item How arXiv tokens influence math-specific tasks beyond the scope of this study, for instance, the process of transforming formal theorems or proofs into more informal descriptions;
    \item The interactions and outcomes resulting from the integration of arXiv tokens with alternative data sources;
    \item If scaling to larger models would amplify or reveal the advantages conferred by incorporating arXiv papers.
\end{itemize}

Thus, further exploration is required, which we leave for future studies.

\subsection{Insights of Reinforcement Learning}
\subsubsection{Towards to a Unified Paradigm} This section proposes a comprehensive framework for examining various training strategies, including SFT, RFT, DPO, PPO, and GRPO. We proceed to experiment in order to investigate key components within this unified framework. Typically, the gradient of a particular training method with regard to the parameter $\theta$ can be expressed as:

\begin{equation}
    \nabla_{\theta}\mathcal{J}_{\textcolor{red}{\mathcal{A}}}(\theta) = \mathbb{E}[\underbrace{(q,o) \sim \textcolor{red}{\mathcal{D}}}_{Data \ Source}]\left( \frac{1}{|o|} \sum_{t=1}^{|o|}  \underbrace{GC_{{\mathcal{A}}}(q, o, t, \textcolor{red}{\pi_{{rf}}})}_{Gradient \ Coefficient}  \nabla_{\theta}\log \pi_{\theta}(o_t | q, o_{<t})\right).
\label{eq:objective}
\end{equation}

There exist three key components: 1) \textit{Data Source $\mathcal{D}$}, which determines the training data; 2) \textit{Reward Function $\pi_{{rf}}$}, which is the source of the training reward signal; 3) \textit{Algorithm $\mathcal{A}$}: which processes the training data and the reward signal to the gradient coefficient $GC$ that determines the magnitude of the penalty or reinforcement for the data. We analyze several representative methods based on such a unified paradigm:

\begin{itemize}
    \item  \textbf{Supervised Fine-tuning (SFT)}: In SFT, the pretrained model is adjusted by training it on a dataset curated by humans for SFT purposes.
    \item \textbf{Rejection Sampling Fine-tuning (RFT)}: RFT takes the model already trained with SFT and further fine-tunes it by using outputs that have been sampled in response to SFT prompts, selecting only those whose answers are deemed correct.
    \item \textbf{Direct Preference Optimization (DPO)}: DPO builds upon the SFT model by fine-tuning it with augmented samples, employing a pair-wise DPO loss function to guide the optimization.
    \item \textbf{Online Rejection Sampling Fine-tuning (Online RFT)}: Unlike traditional RFT, Online RFT starts with the SFT model as the policy model and continually fine-tunes it using new augmented outputs generated from the current policy in real time.
    \item \textbf{PPO/GRPO}: PPO/GRPO commences with the SFT model serving as the policy model and further improves it by reinforcing it on outputs generated online by the current policy model.
\end{itemize}

The elements constituting these approaches are outlined in Table \ref{tab:data_policy}. For a comprehensive exposition of the derivation procedure, see Appendix \ref{app:analysis-rl}.

\paragraph{Observation about Data Source} We classify data sources into two types: online sampling and offline sampling. In online sampling, the training dataset is generated from the exploration outcomes of the currently evolving policy model. In contrast, offline sampling refers to data collected based on the outputs from the original SFT model. Approaches such as RFT and DPO utilize the offline method, whereas Online RFT and GRPO employ the online approach.

As depicted in Figure \ref{fig:rl-analysis}, our analysis reveals that Online RFT consistently yields better performance than RFT across both benchmarks. In particular, while Online RFT performs similarly to RFT at the beginning of training, it demonstrates a clear performance lead as training progresses, highlighting the benefits of online learning methods. This observation aligns with expectations: during the initial training phase, the actor closely mirrors the SFT model, so the variance in sampled data is minimal. However, as training continues, the differences in the actor's samples become more pronounced, allowing online data collection to provide more substantial benefits.

\paragraph{Observation about Gradient Coefficient} In the described approach, the reward signal is transformed into a gradient coefficient, which is then used to modify the model parameters during training. In our experimental design, we partition the reward function into two types: `Rule' and `Model'. The `Rule' component evaluates response quality purely on answer correctness, while the `Model' aspect involves employing a learned reward model to assign scores to each response. The training data used for this reward model is generated by the initial rule-based judgments. As depicted in Equations \ref{eq:GC-RFT} and \ref{eq:GC-GRPO}, a fundamental distinction emerges between GRPO and Online RFT: GRPO singularly tailors the gradient coefficient in response to the value outputted by the reward model, thereby enabling fine-grained rewards and penalties proportional to the magnitude of the responses. Conversely, Online RFT omits this adaptive adjustment; it refrains from penalizing incorrect outputs, instead uniformly rewarding every correct response with the same level of reinforcement.

As illustrated in Figure \ref{fig:rl-analysis}, GRPO outperforms online RFT, underscoring the advantages of adjusting the coefficients for positive and negative gradients. Moreover, GRPO+PS delivers better results than GRPO+OS, demonstrating the value of employing finely tuned, step-aware gradient coefficients. Additionally, our investigation includes iterative RL, with experiments comprising two iterations. As depicted in Figure \ref{fig:iter-rl}, we observe that iterative RL yields notable performance gains, particularly during the initial iteration.

\subsubsection{Why RL Works?} This study applies reinforcement learning using a subset of instruction tuning data and observes that it leads to notable improvements over models trained purely with instruction tuning. To investigate the underlying reasons for the effectiveness of reinforcement learning, we measure both Pass@K and Maj@K accuracy for the Instruct and RL models across two benchmark datasets. Results depicted in Figure \ref{fig:maj-pass} demonstrate that RL substantially boosts Maj@K performance, while Pass@K remains largely unaffected. These outcomes suggest that reinforcement learning chiefly strengthens the reliability of the output distribution, essentially by \textbf{amplifying the likelihood of the correct answer appearing among the TopK results, rather than enhancing the model's core skill set.} In parallel, \citep{wang2023making} observed a \textbf{misalignment problem} in SFT models on reasoning benchmarks, where they show that implementing various preference alignment approaches can noticeably improve reasoning performance \citep{yuan2023rrhf,song2023preference,wang2023making}.

\subsubsection{How to Achieve More Effective RL?}
\label{sec:effec_rl}
Our results indicate that RL is quite successful when applied to tasks involving mathematical reasoning. Additionally, we introduce a coherent framework that encapsulates various prominent training approaches, interpreting each as a form of direct or simplified RL. According to the structure outlined in Equation \ref{eq:objective}, three principal elements emerge: Data Source, Algorithm, and Reward Function. We outline prospective avenues for advancing each of these core components.

\paragraph{Data Source} The foundation of all training approaches is the data source. Within the RL framework, the data source denotes the unlabeled questions along with their corresponding outputs, which are generated by sampling from the policy model. In this work, we rely exclusively on the set of questions introduced during the instruction tuning phase and utilize a straightforward nucleus sampling method for output generation. We hypothesize that this design choice may underlie the observed improvements being limited to Maj@K metrics within our RL pipeline. Looking ahead, we plan to extend our RL pipeline to encompass out-of-distribution question prompts, while also integrating \textbf{advanced sampling (decoding) strategies}, such as those leveraging tree-search algorithms \citep{yao2023tree}. Additionally, \textbf{efficient inference techniques} \citep{xia-etal-2023-speculative,leviathan2023fast,kwon2023efficient,xia2024unlocking}, which significantly influence the efficiency with which policy models explore, will be further investigated due to their critical importance.

\paragraph{Algorithms} Algorithms utilize both data and the reward signal to compute the gradient coefficient, which is then used to update model parameters. As indicated in Equation \ref{eq:objective}, current approaches generally operate under the assumption of complete \textbf{TRUST} in the reward function's feedback, adjusting the conditional probability of generating specific tokens accordingly. Nevertheless, this reward signal cannot always be depended upon, particularly when handling highly intricate tasks. For instance, even the PRM800K datasets \citep{lightman2023let}—meticulously annotated by skilled annotators—still exhibit roughly 20\% incorrect annotations\footnote{\url{https://github.com/openai/prm800k/issues/12\#issuecomment-1728491852}}. Therefore, we focus on reinforcement learning algorithms designed to remain effective despite noisy reward signals. We argue that alignment strategies such as \textbf{WEAK-TO-STRONG} \citep{burns2023weak} possess the potential to fundamentally transform how learning algorithms operate.

\paragraph{Reward Function} The reward function provides the fundamental training signal in RL, and it is typically represented by a neural reward model. We identify three central research avenues for the development of reward models: 1) \textbf{Improving the generalization capacity of reward models.} Effective reward models should generalize across novel, out-of-distribution queries and sophisticated decoded outputs; failure to achieve this may result in reinforcement learning merely reinforcing the prior distribution of LLMs instead of advancing their core abilities; 2) \textbf{Representing the uncertainty inherent in reward models.} Capturing this uncertainty can facilitate the integration of imperfect reward models with learning algorithms designed to leverage weak supervision; 3) \textbf{Constructing high-quality process reward models efficiently,} which can supply detailed and granular training feedback for complex reasoning processes \citep{lightman2023let,wang2023math}.

\section{Conclusion, Limitation, and Future Work} This paper introduces DeepSeekMath, a model that achieves superior results over all open-source counterparts on the MATH competition benchmark and rivals the performance of proprietary systems. DeepSeekMath is built upon DeepSeek-Coder-v1.5 7B and trained further using 500B tokens, with a notable portion—120B tokens—representing mathematical content extracted from Common Crawl. Through comprehensive ablation experiments, we demonstrate that web pages are a valuable resource for curating high-quality mathematics data, whereas arXiv contributes less than initially anticipated. We propose Group Relative Policy Optimization (GRPO), an adaptation of Proximal Policy Optimization (PPO), which enhances mathematical reasoning abilities while reducing memory requirements. Experimental findings confirm that GRPO yields benefits even when DeepSeekMath-Instruct 7B already performs strongly on benchmarks. Additionally, we present a coherent framework for interpreting various methods in this area and outline multiple prospective avenues to develop more efficient reinforcement learning approaches.

While DeepSeekMath~demonstrates strong results on benchmarks for quantitative reasoning, its performance on geometry tasks and theorem proving remains inferior compared to proprietary models. For example, during preliminary testing, the model struggled with questions involving triangles and ellipses, suggesting a potential selection bias in the data used for both pre-training and fine-tuning. Moreover, DeepSeekMath~is constrained by model size, resulting in less robust few-shot learning compared to GPT-4. Unlike GPT-4, which shows notable improvements with few-shot prompts, DeepSeekMath~exhibits comparable outcomes in both zero-shot and few-shot scenarios. To address these issues, future work will focus on refining our data selection processes to build a higher quality pre-training dataset. Furthermore, we plan to investigate additional strategies (see Section \ref{sec:effec_rl}) to enhance the effectiveness of reinforcement learning for LLMs.

\end{document}